\section{Introduction}
An observation can be favorable in two different ways: it can lie near a representable prediction, and it can leave little uncertainty about the protected future test. These are not the same property. In an informative hidden-state experiment, the last report usually does not determine the posterior. Compressing a convenient function of that report while ignoring the remaining conditional variance can therefore give an arbitrarily optimistic resource law.

We study the experiment-level passage
\[
\begin{aligned}
&\text{preparation and positive kernels}\\
&\quad\longrightarrow\text{executable future tests and predictive quotient}\\
&\quad\longrightarrow\text{actual marked acquisition law}\\
&\quad\longrightarrow\text{causal implementation and protected-task risk}.
\end{aligned}
\]
The principal theorem uses a calibrated report proxy $f$ together with a residual-risk mark $w\ge0$. Its assumptions do not assert that $f$ is a sufficient statistic. Instead, at every legal history, the conditional variance above a common irreducible baseline is comparable to $w$, and the squared bias of $f$ is bounded by $w$. These inequalities have an unnormalized one-call formulation in the raw instrument. Consequently they can be checked on an invariant family of prepared physical laws, before any observer or budget is selected.

The resulting geometry is a probability measure on pairs $(f,\sqrt w)$, with distortion $w+\min_j\|f-r_j\|^2$. This is an actual one-call law of a specified prepared reference experiment. Every adaptive stopping procedure induces a terminal submeasure bounded by its paid acquisition occupation. Conversely, a common parameter-independent program uses fresh action and acceptance coins and only one live label to realize the corresponding upper bound. Histories continue to alter the hidden state and the conditional audit mean; no posterior reconstruction from a last report is presumed. Theorem~\ref{thm:central} supplies both directions and an explicit mass--cost dual.

When location and low-uncertainty masses have exponents $d$ and $\alpha$, respectively, the profile has two distinct constraints. Concentrating all terminal mass on small uncertainty costs $N^{-2/\alpha}$. Simultaneously concentrating it near one of $m$ output values costs $(Nm)^{-2/(d+\alpha)}$. Theorem~\ref{thm:geometry} proves that the larger term is both necessary and attainable. For an autonomous observer, $m=M-1$ counts every output label after paying for the live controller. The resulting saturation threshold $m\asymp N^{d/\alpha}$ is a statement about the acquired task geometry, not about the dimension of the physical carrier or the number of stored real coordinates.

Two informative realizations verify this same general theorem. The first carries a genuinely evolving hidden Markov residual through controlled observations. The second retains quantum coherence under noncommuting weak instruments. Both include a reported location and reliability, while the unrevealed task component remains history dependent. Singular regular supports, including Cantor products, are permitted without a smooth manifold assumption.

A separate consequence is a natural Gaussian benchmark. For a hidden sign observed through $N(\pm\mu,1)$ likelihoods, an all-history converse and a three-state construction have the same acquisition rate
\[
 \log R(N)=-2\mu\sqrt{2\log N}+O(1).
\]
The lower bound allows arbitrary stopping and arbitrary terminal readouts; it does not compare only one controller. A revealing branch creates a non-power-law critical transition. In the symmetric classical chain we additionally determine the exact risk of every optimal two-state observer, including randomized rules, and its larger zero-risk budget. These are consequences of the general risk-envelope mechanism and actual occupation balance, not substitutes for the general theorem.

\subsection*{Relation to earlier work}
Conditional projection, randomized selection by conditional risk, and finite-state likelihood-ratio testing are established tools. Selective classification optimizes uncertainty versus coverage~\cite{Franc}; trimmed quantization optimizes a bounded density and a codebook~\cite{CGM}; finite-memory hypothesis testing has a classical theory~\cite{HC}. We do not claim these mechanisms as new. The present comparison concerns a changing protected state, history-dependent report laws, a paid stopping rule, and a complete autonomous state count. Its specific new conclusions are the calibrated marked occupation transfer, the joint reliability--resolution crossover, and the stated informative and singular acquisition laws.

Predictive equivalence~\cite{SC}, value-directed belief compression~\cite{PB}, and Bayes-risk/Bregman geometry~\cite{Varshney} concern related but different information objects. The current finite-memory belief approximation result~\cite{Kim} explicitly compares costs under the same closed-loop execution; our lower bounds quantify over every legal history-dependent stopping policy. Gaussian filtering and weak quantum measurement supply the raw experimental context~\cite{Rabiner,Jacobs}. The retained foundations articles contain the exact common-program occupation criterion, the likelihood-aware three-defect identity, and the report-recoverable acquisition theorem~\cite{GTF49,GTF32,GTF51}. We preserve those results with their original hypotheses. Here the risk mark removes the report-recoverability restriction for the declared calibrated class; it does not claim a necessary and sufficient classification of all causal experiments.
